\section{The comparison with units and ideal classes}
\label{su-section}

This section proves assertion~\ref{sc-count}.  We first give
a proof of the classical Dirichlet unit theorem
\cite[Chapter 5]{milne-ant} from the lattice-volume argument
and ideal count already established.
The comparison below, including the determinant of the action,
the invariant Kummer classes, and all cardinality estimates, is
proved here.  No homology-class construction from
assertion~\ref{sc-few} is assumed except its stated finite-set
bound $\mathcal A(b,B)$.

\subsection{The logarithmic unit lattice}

\begin{theorem}\label{su-unit-theorem}
Let $F$ be a number field of degree $n=r_1+2r_2$.
For one embedding at each infinite place put
\[
 \lambda(u)=(\log|\sigma_vu|)_v,\qquad
 H=\{x:\ \sum_v w_vx_v=0\},
 \quad
 w_v=\begin{cases}1&v\text{ real},\\2&v\text{ complex}.\end{cases}
\]
The image $\lambda(\mathcal O_F^\times)$ is a full lattice
in $H$, and its kernel is the finite cyclic group $\mu(F)$.
Consequently $\mathcal O_F^\times/\mu(F)$ is free abelian
of rank $r_1+r_2-1$.
\end{theorem}
\begin{proof}
A unit has norm of absolute value one, so its logarithm
belongs to $H$.  If these logarithms belong to a bounded
set, all conjugates of the units are bounded.
By (SL1), only finitely many such algebraic integers
exist.  Thus the logarithmic image has finite intersection
with every bounded subset of $H$, and is discrete.
Its kernel is a finite group.  Every element of this
finite group is a root of unity, and every root of unity
lies in the kernel.  A finite subgroup of
$\mathbb C^\times$ is cyclic: if its exponent is $m$,
the equation $X^m=1$ has at most $m$ roots, whereas the
classification of finite abelian groups gives at least
$m$ elements, with equality only for a cyclic group.
Equivalently, it is a subgroup of the cyclic group of
complex $m$-th roots of unity.

It remains to prove that the logarithmic image spans $H$.
Choose $C>1$ so large that
\[
       2^{r_1}\pi^{r_2}C^n
                         >2^n2^{-r_2}\sqrt{D_F}.
\]
For each $x\in H$, apply the lattice-volume lemma to the
product of intervals and disks
\[
    |y_v|<Ce^{x_v}\quad(v\text{ real}),\qquad
    |z_v|<Ce^{x_v}\quad(v\text{ complex}).
\]
Its volume is $2^{r_1}\pi^{r_2}C^n$, independent of $x$,
because $\sum_vw_vx_v=0$.  Hence it contains
$0\ne\alpha\in\mathcal O_F$.
Write $b_v=\log|\sigma_v\alpha|-x_v$.
Each $b_v<\log C$, and
\[
          0\le\log|N_{F/\mathbb Q}\alpha|
                    =\sum_vw_vb_v<n\log C.
\]
The lower inequality holds because the nonzero norm is
an integer.  For each $v$ it follows that
\[
          -\frac{n-w_v}{w_v}\log C
                    \le b_v<\log C .
\]
Thus $\lambda(\alpha)-x$ stays in a fixed bounded box.

The principal ideal $(\alpha)$ has norm at most $C^n$.
By Lemma~\ref{sl-ideal-count}, only finitely many such
ideals occur.  Choose a generator $\beta_i$ for each
principal ideal in that finite list.  If
$(\alpha)=(\beta_i)$, then $u=\alpha/\beta_i$ is a unit,
and
$\lambda(u)-x=\lambda(\alpha)-x-\lambda(\beta_i)$
lies in a bounded set independent of $x$.
Consequently every point of $H$ is a bounded distance
from a logarithm of a unit.  A proper real linear
subspace of $H$ cannot have this property, as points
on a line transverse to it have unbounded distance.
The logarithmic image therefore spans $H$.

Choose a real basis of $H$ from this discrete subgroup.
The half-open parallelepiped of its integral span is
bounded, so it contains only finitely many elements of
the subgroup.  Reduction into that parallelepiped shows
that the integral span has finite index in the subgroup.
The subgroup is thus finitely generated, torsion-free,
and of rank $\dim H=r_1+r_2-1$; its parallelepipeds
give a lattice in all of $H$.  This proves the theorem,
including $\dim H=0$.
\end{proof}

\subsection{The divisor sequence and the finite comparison group}

Fix the data (SC3), with $K/L$ cyclic of degree
$d=p^{\nu-c}$, and put
\[
 q=p^\nu,\qquad a=(1+p^c)^E,\qquad
 v=v_p(E),\qquad r=r_2(L).
\]
Thus $a>1$, $a\equiv1\pmod p$, and $L$ is totally imaginary.
The hypotheses also imply $K=L(\mu_{p^\nu})$:
the restriction of $\gamma$ to the primitive root has order
$p^{\nu-c}=d$ by (SC18), so it has trivial kernel in
$\operatorname{Gal}(K/L)$.
We write $Y=K^\times$ additively for the next group calculations:
$y+y'$ means the product of the corresponding nonzero field
elements, and $my$ means their $m$-th power.
Let $I=\bigoplus_P\mathbb Z P$ be the group of fractional
divisors and let $\partial:Y\to I$ be the principal-divisor
map.  Its kernel is the unit group $U_K=\mathcal O_K^\times$,
and its cokernel is the ideal class group $C_K$.
Write $\mu=\mu(K)$ for the torsion subgroup of $Y$.

Let $g$ act on $Y$ by $\gamma^{-1}$ and on $I$ by
$g(P)=\gamma^{-1}P$.  In coordinates,
$(gD)_P=D_{\gamma P}$, so $\partial$ commutes with $g$.
Set
\[
 \delta=g-a,\qquad
 V=\{y\in Y:\partial y\in\delta I\},\qquad
 W=V/(\delta Y+\mu),\qquad
 C=C_K/\delta C_K .
 \tag{SU1}
\]
Every torsion element has zero divisor, and
$\partial(\delta Y)\subset\delta I$, so this quotient is defined.
Multiplicatively, $y\in V$ means exactly that it is admissible
in (SC4).  Since $g^d=1$ and $a>1$, $\delta$ is injective on
the torsion-free group $I$.  Indeed $\delta x=0$ implies
$x=a^dx$, and $a^d-1\ne0$.

\begin{lemma}\label{su-exact}
There is an exact sequence of finite groups
\[
 0\longrightarrow
 (U_K/\mu)/\delta(U_K/\mu)
 \longrightarrow W
 \longrightarrow\ker(\delta:C_K\to C_K)
 \longrightarrow0.
 \tag{SU2}
\]
Moreover
\[
 |W|=\frac{(a^d-1)^r}{a-1}\,|C|,
 \qquad
 v_p\!\left(\frac{(a^d-1)^r}{a-1}\right)
                    =r(\nu+v)-c-v .
 \tag{SU3}
\]
\end{lemma}
\begin{proof}
For $y\in V$ there is a unique $D\in I$ with
$\delta D=\partial y$.  Send the class of $y$ to the ideal
class $[D]$.  Changing $y$ by $\delta z+t$, $t\in\mu$,
changes $D$ by $\partial z$, so the map is well-defined.
Its image is precisely $\ker\delta$ on $C_K$: an ideal
class killed by $\delta$ has a representative $D$ with
$\delta D$ principal.  Its kernel is represented by units,
because $D=\partial z$ implies $y-\delta z\in U_K$.
Finally
\[
             U_K\cap(\delta Y+\mu)=\delta U_K+\mu:
\]
if a unit is $\delta z+t$, taking divisors and using
injectivity of $\delta$ on $I$ gives $\partial z=0$.
This proves the exact sequence.

Theorem~\ref{su-unit-theorem} identifies $(U_K/\mu)\otimes\mathbb R$
with the hyperplane $H$ of coordinate sum zero in
$\mathbb R^{r_2(K)}$.  All its places are complex, so
the weights in the product formula are equal.
The $d$ automorphisms of $K/L$ act freely on the complex
places over each of the $r$ complex places of $L$.
To verify freeness, fix one complex embedding of $L$.
Its $d$ extensions to $K$ form a free Galois orbit;
none is identified with another by conjugation, since
that would change the chosen embedding of $L$.
Thus the permutation of coordinates is a disjoint union
of $r$ cycles of length $d$.

The characteristic polynomial on each cycle is $X^d-1$.
The quotient of $\mathbb R^{rd}$ by $H$ is the trivial
one-dimensional representation.  Therefore
\[
        |\det(\delta|_H)|=\frac{(a^d-1)^r}{a-1}.
\]
An injective endomorphism of a free integral lattice has
cokernel order equal to the absolute determinant of its
matrix: integer row and column operations reduce the
matrix to Smith normal form.  It follows that the unit
quotient is finite with the displayed order.
The group $C_K$ is finite by Lemma~\ref{sl-small-ideal}.
For an endomorphism of a finite group, the kernel and
cokernel have the same order; both orders are the total
order divided by the image order.  This proves the first
formula (SU3) and finiteness of $W$.

Finally (SC18) gives
$v_p(a-1)=c+v$ and $v_p(a^d-1)=\nu+v$.
Taking valuations proves the second formula.
\end{proof}

The equivariant class-group functions in (SC6) identify with
$\operatorname{Hom}(C,\mathbb Z/q\mathbb Z)$.
Indeed vanishing on principal divisors factors a function
through $C_K$, and vanishing on $(g-a)P$ is precisely
$\varphi(\gamma^{-1}P)=a\varphi(P)$.
For a finite abelian group $A$, write $|A|_p=p^{v_p(|A|)}$.
Decomposition into cyclic groups gives
\[
 |\operatorname{Hom}(A,\mathbb Z/p^\nu\mathbb Z)|
        =\prod_i p^{\min(e_i,\nu)}\le |A|_p,
 \tag{SU4}
\]
where $p^{e_i}$ are the orders of its cyclic $p$-primary
factors.  Consequently (SU3) implies
\[
            |T_{\gamma,a,q}|q^r
                   \le |W|_p\,p^{c+v}.
                   \tag{SU5}
\]
In obtaining this inequality we discarded the additional
nonnegative exponent $rv$ in the unit factor.

\subsection{A uniform bound for the \texorpdfstring{$p$}{p}-rank}

Define
\[
 \begin{split}
 \mathcal K_p(L)
   &=\{x\in L^\times/(L^\times)^p:
                 v_{\mathfrak p}(x)\equiv0\pmod p
                       \text{ if }\mathfrak p\nmid p\},\\
 \mathcal I_p(K)
   &=\{[y]\in Y/(pY+\mu):
            gy-y\in pY+\mu,\quad \partial y\in pI\}.
 \end{split}
 \tag{SU6}
\]
These definitions are independent of representatives.
The group $\mathcal K_p(L)$ is a Kummer group with the
displayed valuation conditions, not the strict Selmer group.

\begin{lemma}\label{su-base-kummer}
If $n=[L:\mathbb Q]$ and $h_L$ is the class number, then
\[
           |\mathcal K_p(L)|
                \le p^{\,2n+\lfloor\log_2h_L\rfloor}.
                \tag{SU7}
\]
\end{lemma}
\begin{proof}
First impose the valuation condition also at primes above
$p$, and call the resulting subgroup $\mathcal K_p^0(L)$.
Sending $\alpha$ to the ideal class of
$\frac1p\operatorname{div}(\alpha)$ maps this group to
$C_L[p]$.  Its kernel consists exactly of the image of
$\mathcal O_L^\times/(\mathcal O_L^\times)^p$:
if the divided divisor is principal, dividing $\alpha$
by the corresponding $p$-th power gives a unit.
By Theorem~\ref{su-unit-theorem} and cyclicity of roots of unity this
kernel has order at most $p^{r_1(L)+r_2(L)}\le p^n$.
If $|C_L[p]|=p^t$, then $2^t\le p^t\le h_L$, whence
$t\le\lfloor\log_2h_L\rfloor$.

Reduction of the remaining valuations modulo $p$ injects
$\mathcal K_p(L)/\mathcal K_p^0(L)$ into one copy of
$\mathbb Z/p\mathbb Z$ for each prime above $p$.
There are at most $n$ such primes.
Multiplying the three bounds proves finiteness and (SU7).
\end{proof}

\begin{lemma}[Cyclic Hilbert theorem 90]\label{su-hilbert90}
If $K/L$ is cyclic with generator $\gamma$ and
$w\in K^\times$ has norm one, then $w=\gamma(z)/z$
for some $z\in K^\times$.
\end{lemma}
\begin{proof}
Put $c_0=1$ and
$c_j=(w\gamma(w)\cdots\gamma^{j-1}(w))^{-1}$ for $j\ge1$.
For $b\in K$ define $z_b=\sum_{j=0}^{d-1}c_j\gamma^j(b)$.
The norm condition gives $c_d=1$, and shifting the summation
index gives $\gamma(z_b)=wz_b$.
For some $b$, $z_b\ne0$, by linear independence of distinct
field automorphisms over $K$.
Here is the usual short proof of that independence.
A nonzero relation with the fewest terms, after rescaling,
has coefficient one at its first automorphism.
Choose $x$ on which a second automorphism differs from the
first, evaluate the relation at $xy$, and subtract the
relation at $y$ multiplied by the first automorphism of $x$.
This is a nonzero relation with fewer terms, a contradiction.
A single automorphism cannot have a nonzero scalar multiple
equal to the zero map.  Taking $z=z_b$ proves the lemma.
\end{proof}

\begin{lemma}\label{su-invariant-kummer}
For the cyclotomic data under consideration,
\[
                |\mathcal I_p(K)|
                     \le p^2|\mathcal K_p(L)|.
                     \tag{SU8}
\]
\end{lemma}
\begin{proof}
Choose a generator $\rho$ of the finite cyclic group
$\mu(K)$.  Its finiteness also follows directly from
(SL1): roots of unity are integral and all their conjugates
have absolute value one, so they belong to a compact set
in a lattice.  A finite subgroup of $\mathbb C^\times$
is cyclic, and any complex embedding proves cyclicity here.

Every class in $\mathcal I_p(K)$ has a representative
$\alpha\in K^\times$ with all valuations divisible by $p$
and with
\[
           \gamma(\alpha)=\alpha w^p\rho^j,
              \qquad 0\le j<p,\quad w\in K^\times .
              \tag{SU9}
\]
The invariance condition with $\gamma^{-1}$ is equivalent
to that with $\gamma$.  Reduction of the exponent of
$\rho$ modulo $p$, absorbing its $p$-th power into $w$,
proves (SU9).
Taking ordinary field norms gives
\[
                     (N_{K/L}w)^p=(N_{K/L}\rho)^{-j}.
\]
There are at most $p$ choices for $j$ and at most $p$
values of $N_{K/L}w$ for each $j$, since a polynomial of
degree $p$ over a field has at most $p$ roots.

Compare two representatives $\alpha,\alpha'$ with the
same $j$ and the same norm of $w,w'$.  Put
$\beta=\alpha/\alpha'$ and $v=w/w'$.
Then $\gamma(\beta)/\beta=v^p$ and $N_{K/L}v=1$.
By Lemma~\ref{su-hilbert90}, $v=\gamma(z)/z$.
It follows that $\beta=xz^p$ with $x\in L^\times$.
The extension $K=L(\mu_{p^\nu})$ is unramified outside $p$:
the proof of Lemma~\ref{sl-cyclotomic-discriminant} puts
$p^\nu$ in its relative different, so that the different
is the unit ideal at every other prime.
Thus for $P\mid\mathfrak p\nmid p$,
$v_P(x)=v_{\mathfrak p}(x)$.
The valuations of $\alpha$ and $\alpha'$ and those of
$z^p$ show that this number is divisible by $p$.
Hence $[x]\in\mathcal K_p(L)$.

For each of the at most $p^2$ pairs
$(j,N_{K/L}w)$ that occurs, fix one representative.
Every class with that pair differs from its representative
by the image of a class in $\mathcal K_p(L)$.
This covers $\mathcal I_p(K)$ by at most $p^2$ sets of
size at most $|\mathcal K_p(L)|$, proving (SU8).
No uniqueness of the pair assigned to a class is needed.
\end{proof}

\begin{lemma}\label{su-rank}
With $g_L$ as in (SC1),
\[
                           |W[p]|\le p^{g_L}.
                           \tag{SU10}
\]
\end{lemma}
\begin{proof}
There is an injection $W[p]\to\mathcal I_p(K)$.
To construct it, choose $u\in V$ representing a class
$w\in W[p]$, and write
\[
                         pu=\delta y+t,\qquad t\in\mu .
                         \tag{SU11}
\]
Such $y,t$ exist by the definition of $W$.
Write $\partial u=\delta D$.
Taking divisors in (SU11) and using injectivity of
$\delta$ on $I$ gives $\partial y=pD$.
Since $p\mid a-1$, equation (SU11) also gives
$gy-y\in pY+\mu$.  Thus $[y]\in\mathcal I_p(K)$.

This class is independent of choices.  If $u$ changes by
$\delta z+t_0$, then subtraction of the two equations
(SU11) shows
$\delta(y'-y-pz)\in\mu$.  Whenever $\delta x\in\mu$,
one has $x\in\mu$: apply the polynomial
$g^{d-1}+ag^{d-2}+\cdots+a^{d-1}$ and use
$g^d=1$ to obtain $(1-a^d)x\in\mu$.
The torsion subgroup of an abelian group is saturated,
meaning that a nonzero multiple of an element is torsion
only if the element is torsion.  Therefore
$y'-y-pz\in\mu$, as required.
If $[y]=0$, write $y=pz+t_1$; then (SU11) gives
$p(u-\delta z)\in\mu$, so $u-\delta z\in\mu$ and $w=0$.
The construction respects addition, hence it is injective.
Combine (SU7), (SU8), and
$g_L=2n+\lfloor\log_2h_L\rfloor+2$.
\end{proof}

\subsection{The weighted product and its kernel}

In the additive notation set
\[
       N_a=\sum_{j=0}^{d-1}a^j g^{d-j},\qquad
       s=p^b,\qquad
       R=\frac{a^d-1}{q}.
       \tag{SU12}
\]
The symbol $R$ denotes an integer in this subsection.
Since $g=\gamma^{-1}$ and $\gamma^d=1$, the element
$N_a y$ is exactly the weighted product in (SC5).
Telescoping gives
\[
                    \delta N_a=N_a\delta=(1-a^d)g=-qRg.
                    \tag{SU13}
\]
Also $v_p(R)=v$, by (SC18).

\begin{lemma}\label{su-weighted-map}
The formula
\[
          f:W\longrightarrow Y/(qY+\mu),\qquad
                     [y]\longmapsto[sN_a y]
                     \tag{SU14}
\]
defines a homomorphism.  Its kernel is killed by $sR$.
If $\mathcal A(b,B)$ holds, then $|\operatorname{im}f|\le B$.
\end{lemma}
\begin{proof}
Equation (SU13) shows that $sN_a(\delta Y)\subset qY$.
An endomorphism of an abelian group preserves its torsion
subgroup, so the formula descends to $W$.
If $f([y])=0$, write $sN_a y=qz+t$ with $t\in\mu$.
Apply $\delta$ and use (SU13):
\[
                      q(-sRg y-\delta z)=\delta t.
\]
The right-hand side is torsion.  Saturation of $\mu$
gives $sRg y\in\delta Y+\mu$.
The automorphism $g$ preserves all groups in (SU1);
applying its inverse shows $sRy\in\delta Y+\mu$.
Thus $sR$ kills the kernel.

The group $W$ is finite by Lemma~\ref{su-exact}.
Choose one representative $y\in V$ for each value of $f$.
These give distinct elements $\alpha\in K^\times$, each
admissible for a divisor.  If two weighted products raised
to $p^b$ differ by a $q$-th power, their images in the
even coarser quotient $Y/(qY+\mu)$ coincide.
They were chosen for distinct values, so their
representatives are equal.  This finite set therefore
satisfies exactly (SC7).  Assertion $\mathcal A(b,B)$
bounds its size by $B$.
We have not replaced equality modulo powers by equality
modulo roots of unity in the hypothesis.
\end{proof}

\begin{lemma}\label{su-finite-group}
Let $W$ be a finite abelian group and $f:W\to A$ a
homomorphism.  Suppose $|W[p]|\le p^g$ and the
$p$-primary part of $\ker f$ is killed by $p^k$.
Then
\[
                         |W|_p\le|\operatorname{im}f|p^{kg}.
                         \tag{SU15}
\]
\end{lemma}
\begin{proof}
Multiplication by $p$ maps $W[p^{j+1}]$ into
$W[p^j]$ with kernel $W[p]$.  Induction gives
$|W[p^k]|\le|W[p]|^k\le p^{kg}$.
The $p$-primary part of $\ker f$ is a subgroup of
$W[p^k]$.  Now use $|W|=|\ker f||\operatorname{im}f|$
and take $p$-primary parts, bounding that of the image
by its full order.
\end{proof}

\begin{theorem}\label{su-count}
Assertion~\ref{sc-count} holds.  More explicitly,
\[
 |T_{\gamma,a,q}|q^{r_2(L)}
       \le B\,p^{(b+v_p(E))g_L+c+v_p(E)}
       \quad\text{if }\mathcal A(b,B)\text{ holds}.
 \tag{SU16}
\]
\end{theorem}
\begin{proof}
Write $R=p^vR_0$ with $p\nmid R_0$.
Lemma~\ref{su-weighted-map} shows that the $p$-primary
kernel of $f$ is killed by $p^{b+v}$, because
multiplication by $R_0$ is invertible on a finite
$p$-group.
Lemmas~\ref{su-rank} and \ref{su-finite-group} give
$|W|_p\le Bp^{(b+v)g_L}$.
Substitution in (SU5) proves the theorem.
\end{proof}

The factor $q^{r_2(L)}$ has thus been derived from the
determinant of the unit action, rather than inserted as
a rank heuristic.  The extra two powers of $p$ in $g_L$
come from the root-of-unity exponent and the norm value
in (SU9).  The argument covers the boundary case
$\nu=c$, when $K=L$ and $d=1$, and it retains the
condition $c\ge2$ at $p=2$ needed in (SC18).


